%*******************************************************
% Abstract
%*******************************************************
%\renewcommand{\abstractname}{Abstract}
\pdfbookmark[1]{Abstract}{Abstract}
% \addcontentsline{toc}{chapter}{\tocEntry{Abstract}}
\begingroup
\let\clearpage\relax
\let\cleardoublepage\relax
\let\cleardoublepage\relax

\chapter*{Abstract}

Ultracold mass-imbalanced Fermi-Fermi mixtures are predicted to host exotic many-body phases, including superfluid states analogous to those conjectured in unconventional superconductors, under experimentally more accessible conditions than their mass-balanced counterparts. These systems, with their exceptional tunability of interactions via magnetic Feshbach resonances, offer a uniquely controllable platform to explore these phenomena. This thesis presents a thorough experimental characterization of the interaction properties and Feshbach molecules of a mass-imbalanced Fermi-Fermi mixture of $^{161}$Dy and $^{40}$K, a system combining a strongly magnetic lanthanide with an alkali atom, for which no theoretical prediction of the interaction spectrum is available.

In the first part of the thesis, we report on the characterization 
of the interspecies interaction spectrum. A thermalization-based 
Feshbach spectroscopy identifies more than $40$ interspecies 
resonances over a wide range of magnetic fields. Among these, we 
focus on the low-magnetic-field region below \SI{10}{G}, where 
five resonances are identified and characterized through 
measurements of the molecular binding energy. A particularly 
well-isolated resonance near \SI{7}{G}, exhibiting a meaningful 
universal regime, is identified as the most suitable candidate for 
interaction tuning in the strongly interacting regime. We 
characterize interactions at this resonance in detail through the 
study of collective dipole-mode oscillations, observing the 
collisional hydrodynamic crossover and reconstructing the full 
dipole-mode spectrum as a function of scattering length. The 
measurements additionally allow us to characterize the interspecies 
drag effect through the study of a purely dissipative oscillation 
mode.

The second part of the thesis is dedicated to the study of DyK Feshbach molecules associated at the \SI{7}{G} resonance. We develop a protocol to associate and purify an optically trapped sample of DyK dimers, and independently characterize the resonance parameters through two complementary measurements of the molecular bound state. We then perform a systematic study of the dominant loss processes limiting the molecular lifetime. Light-induced losses, driven by photons from the optical dipole trap, are characterized at four different trapping wavelengths between \SI{1000}{nm} and \SI{2000}{nm}, revealing a strong wavelength dependence and identifying conditions under which losses are substantially suppressed. In a trap operating near \SI{1550}{nm}, where light-induced losses are reduced, we observe inelastic collisional contributions to the molecular decay and demonstrate Pauli suppression of two-body losses for very weakly bound dimers.%, constituting direct evidence of the fermionic composite nature of the molecules near the resonance pole.


Together, these results establish the \SI{7}{G} resonance as a 
well-characterized platform for studying strongly interacting 
physics in the Dy-K mixture, and represent a significant step 
toward the realization of a molecular Bose-Einstein condensate of 
DyK dimers and the exploration of the BEC-BCS crossover in a 
mass-imbalanced Fermi system.


\newpage

\begin{otherlanguage}{ngerman}
\pdfbookmark[1]{Zusammenfassung}{Zusammenfassung}
\chapter*{Zusammenfassung}
Kurze Zusammenfassung des Inhaltes in deutscher Sprache\dots
\end{otherlanguage}

\endgroup

\vfill
